\documentclass[11pt,a4paper]{article}
\usepackage[utf8]{inputenc}
\usepackage[T1]{fontenc}
\usepackage[french]{babel}
\usepackage[top=2cm,bottom=2cm,left=2cm,right=2cm]{geometry}
\usepackage{xcolor}
\usepackage{tcolorbox}
\tcbuselibrary{skins,breakable}
\usepackage{fancyhdr}
\usepackage{enumitem}
\usepackage{lmodern}

% --- Couleur discipline : MANAGEMENT ---
\definecolor{mainColor}{RGB}{0,120,60}
\definecolor{darkGray}{RGB}{50,50,50}

\tcbset{boxrule=0.5pt, arc=3pt, left=5pt, right=5pt, top=3pt, bottom=3pt, breakable}

\newtcolorbox{defbox}[1]{
  colback=mainColor!8, colframe=mainColor, coltitle=white,
  fonttitle=\bfseries\footnotesize, title=DÉFINITION \textemdash\ #1}

\newtcolorbox{keybox}{
  colback=darkGray!7, colframe=darkGray!50, coltitle=white,
  fonttitle=\bfseries\footnotesize, title=POINT CLÉ}

\newtcolorbox{exbox}{
  colback=mainColor!5, colframe=mainColor!40,
  fonttitle=\bfseries\footnotesize, title=EXEMPLE}

\pagestyle{fancy}\fancyhf{}
\fancyhead[L]{\small\textcolor{mainColor}{\textbf{CEJM — BTS SIO}}}
\fancyhead[C]{\small\textbf{Chapitre 5 — Finalités et parties prenantes}}
\fancyhead[R]{\small\textcolor{mainColor}{\textbf{MANAGEMENT}}}
\fancyfoot[C]{\small\thepage}
\renewcommand{\headrulewidth}{0.5pt}
\setlength{\headheight}{15pt}
\setlist{itemsep=2pt, topsep=3pt, parsep=0pt}

\begin{document}

\begin{center}
  {\Large\bfseries\textcolor{mainColor}{Chapitre 5}}\\[2pt]
  {\large Les finalités de l'entreprise et les parties prenantes}\\[4pt]
  \textit{\textcolor{darkGray}{Comment concilier performance économique et responsabilité sociétale ?}}
\end{center}
\vspace{4pt}\hrule\vspace{8pt}

\section*{Définitions essentielles}

\begin{defbox}{Finalité}
La \textbf{raison d'être} de l'entreprise — mission fondamentale qui justifie son existence sur le long terme.
\end{defbox}

\vspace{4pt}
\begin{defbox}{RSE — Responsabilité Sociétale des Entreprises}
Intégration \textbf{volontaire} par les entreprises de préoccupations \textbf{sociales, environnementales et économiques} dans leur activité et leurs relations avec les parties prenantes.
\end{defbox}

\vspace{4pt}
\begin{defbox}{Parties prenantes (Stakeholders)}
Ensemble des acteurs \textbf{internes et externes} qui peuvent affecter ou être affectés par les décisions de l'entreprise.
\end{defbox}

\vspace{4pt}
\begin{defbox}{Performance globale}
Résultat de l'entreprise mesuré sur \textbf{trois dimensions} : économique, sociale et environnementale (Triple Bottom Line).
\end{defbox}

\section*{Les 4 finalités de l'entreprise}

\begin{center}
\begin{tabular}{|l|l|}
\hline
\textbf{Finalité} & \textbf{Objectif} \\
\hline
\textbf{Économique} & Produire, être rentable, assurer la pérennité \\
\textbf{Sociale} & Bien-être des salariés, emploi, formation, égalité \\
\textbf{Environnementale} & Réduire l'impact sur l'environnement (CO2, déchets…) \\
\textbf{Sociétale} & Contribuer positivement à la société, éthique des affaires \\
\hline
\end{tabular}
\end{center}

\begin{exbox}
\textbf{Veja (baskets)} : économique (vente et rentabilité), sociale (conditions de travail, insertion handicap), environnementale (coton agroécologique), sociétale (juste rémunération des producteurs).
\end{exbox}

\section*{Les parties prenantes}

\subsection*{Parties prenantes internes}
\begin{itemize}
  \item \textbf{Salariés} : conditions de travail, rémunération, sens au travail
  \item \textbf{Dirigeants} : rentabilité, pérennité, stratégie
  \item \textbf{Actionnaires} : dividendes, valeur de l'action, risques extra-financiers
\end{itemize}

\subsection*{Parties prenantes externes}
\begin{itemize}
  \item \textbf{Clients} : qualité, prix, éthique, transparence
  \item \textbf{Fournisseurs} : relations équitables et durables
  \item \textbf{État} : respect de la réglementation, emploi, impôts
  \item \textbf{Société civile / ONG} : impact environnemental et social
\end{itemize}

\begin{keybox}
L'entreprise doit \textbf{équilibrer} les attentes parfois contradictoires de ses parties prenantes.
\end{keybox}

\section*{La démarche RSE}

\subsection*{Les bénéfices d'une démarche RSE}
\begin{itemize}
  \item \textbf{Image et réputation} : attractivité auprès des clients et candidats
  \item \textbf{Fidélisation} : rétention des talents (marque employeur)
  \item \textbf{Innovation} : nouvelles solutions durables = nouveaux marchés
  \item \textbf{Gestion des risques} : anticipation des risques réglementaires
  \item \textbf{Économies} : réduction des coûts énergétiques et de matières
\end{itemize}

\section*{La performance globale — indicateurs}

\begin{center}
\begin{tabular}{|l|l|}
\hline
\textbf{Dimension} & \textbf{Indicateurs} \\
\hline
Économique & CA, résultat net, rentabilité \\
Sociale & Turn-over, taux d'absentéisme, index égalité H/F \\
Environnementale & Bilan carbone, part d'énergies renouvelables \\
Sociétale & Mécénat, achats solidaires, score RSE \\
\hline
\end{tabular}
\end{center}

\section*{Points clés à retenir}

\begin{keybox}
L'entreprise ne poursuit pas uniquement un but lucratif : ses \textbf{finalités} orientent toute sa stratégie.
\end{keybox}
\vspace{3pt}
\begin{keybox}
La \textbf{RSE} est la mise en œuvre concrète des finalités sociales, environnementales et sociétales.
\end{keybox}
\vspace{3pt}
\begin{keybox}
La \textbf{performance globale} (triple bilan) dépasse la simple rentabilité financière.
\end{keybox}

\section*{Mots-clés}
\textcolor{mainColor}{\textbf{Finalité}} $\cdot$
\textcolor{mainColor}{\textbf{RSE}} $\cdot$
\textcolor{mainColor}{\textbf{Parties prenantes}} $\cdot$
\textcolor{mainColor}{\textbf{Performance globale}} $\cdot$
\textcolor{mainColor}{\textbf{Triple Bottom Line}} $\cdot$
\textcolor{mainColor}{\textbf{Développement durable}} $\cdot$
\textcolor{mainColor}{\textbf{Marque employeur}} $\cdot$
\textcolor{mainColor}{\textbf{Éthique des affaires}} $\cdot$
\textcolor{mainColor}{\textbf{Bilan carbone}}

\end{document}
